\section{Ethical and participatory checklist}
\label{supp:ethics}

This supplementary section is a working checklist for researchers who
intend to use GAWorld v2 in a study that touches identifiable people,
sensitive attributes, or consequential applications.  It is not a
replacement for institutional review; it is a reminder of the
specific risks that an LLM-based artificial-society platform
introduces and of the mitigations that the platform already provides.

\subsection{Identifiable profiles and personas}

A \emph{persona distillation} takes a name or a URL and produces a
sourced portrait that can be deployed as a resident.  The portrait
is held under \texttt{output/personas/<slug>/} and only an explicit
\emph{deploy} step writes the seed CSV that a researcher can publish.
A researcher who intends to use a portrait derived from a real
person should:

\begin{itemize}[leftmargin=*]
  \item \textbf{Document the source} (URL, archive, date of access)
  in the portrait's \texttt{research.md};
  \item \textbf{Document the consent} or the public-figure
  justification for distillation, and the institution's policy on
  distillation;
  \item \textbf{Separate the distill from the deploy}: the portrait
  may be inspected and reviewed without ever producing a resident;
  \item \textbf{Audit the deployed resident} against the portrait and
  against the source; record the SHA-256 of both;
  \item \textbf{Restrict access} to the deployed resident; do not
  include it in a population whose output is public.
\end{itemize}

\subsection{Sensitive attributes and within-group flattening}

Language models may reproduce stereotypes and suppress within-group
heterogeneity.  v2 provides a personality-anchor rewriter that
calibrates Big-Five proxies against a corpus and warns when one
trait dominates a group that the configuration does not intend to
dominate.  A researcher who uses a sensitive attribute (race,
gender, religion, disability, sexual orientation, immigration
status) should:

\begin{itemize}[leftmargin=*]
  \item \textbf{Document the attribute source} (census, registry,
  self-report, model generation);
  \item \textbf{Document the within-group variability} the
  configuration is intended to represent;
  \item \textbf{Run the personality-anchor rewriter} and report
  the resulting distribution;
  \item \textbf{Audit outputs} for stereotype markers; do not
  report aggregate outputs as evidence of population behavior
  without a within-group breakdown.
\end{itemize}

\subsection{Consequential applications}

A consequential application is one whose output can affect a
person's housing, health, employment, policing, mobility, or
information governance.  GAWorld v2 must not substitute for
engagement with people affected by a research question or policy.
A researcher who intends to use v2 in a consequential study should:

\begin{itemize}[leftmargin=*]
  \item \textbf{Engage affected communities} in problem
  formulation, outcome choice, and release conditions;
  \item \textbf{Document the engagement} in a process record that
  survives the study;
  \item \textbf{Map simulated constructs to external measures}
  before interpreting an aggregate output;
  \item \textbf{Reserve targets} for testing rather than
  calibration; do not tune to a target and then report agreement;
  \item \textbf{Report uncertainty} across seeds, plausible
  parameterizations, and counterfactual pipelines;
  \item \textbf{Do not release} a model, a population, or an
  archive that could re-identify a real person.
\end{itemize}

\subsection{Human participation and player interventions}

A player intervention is a real, inspectable input to a simulation;
it is not a ground truth.  A researcher who uses an
\emph{open} world with player interventions should:

\begin{itemize}[leftmargin=*]
  \item \textbf{Document the participant pool} and the recruitment
  procedure;
  \item \textbf{Document the consent} and the data retention
  policy;
  \item \textbf{Separate player interventions from automatic
  interventions} in the analysis; do not pool them;
  \item \textbf{Do not infer} that a player's claim is faithful to
  the claimed resident; report it as a real intervention, not as
  evidence about the resident.
\end{itemize}

\subsection{Audit trail and reproducibility}

Every mutation that crosses the v2 API is recorded in the
\texttt{audit} table, with the caller's user id, the world's id, the
operation type, the target entity, and a timestamp.  A researcher
who intends to publish a study should:

\begin{itemize}[leftmargin=*]
  \item \textbf{Retain} the \texttt{audit} table for the period
  required by the institution;
  \item \textbf{Bind} every reported number to a hash, an
  evidence class, the producing code path, and the world and node
  identifiers;
  \item \textbf{Disclose} the provider, the model, the prompt
  version, and the sampling parameters used in each run;
  \item \textbf{Estimate} replay variability rather than promise
  bitwise determinism for hosted models;
  \item \textbf{Submit} the artifact set to a stable repository
  (Zenodo, Dryad, or the institution's archive) before
  publication.
\end{itemize}

\subsection{Withdrawal and deletion}

A participant or a deployed profile can be withdrawn from a study
by removing the world's seed directory, the run artifacts, and the
audit rows that reference the participant.  v2 provides no
automatic redaction; the redaction is a researcher-driven
operation.  A withdrawal request should be honored within the
period required by the institution and the consent agreement.
